\uuid{T2Ss}
\titre{Intervalle centré pour une loi normale}
\auteur{Maxime Nguyen}
\organisation{AMSCC}
\datecreate{2026-11-24}
\niveau{L2}
\module{Probabilité}
\chapitre{Probabilité continue}
\sousChapitre{Loi normale}
\difficulte{2}
\theme{loi normale, standardisation, intervalle centré}

\contenu{
\begin{question}
\qtype{onlyone}
On note $\mathcal{N}(\mu,\sigma)$ la loi normale d'espérance $\mu$ et d'écart-type $\sigma$, et $\Phi$ la fonction de répartition de $\mathcal{N}(0,1)$. Soit $X \sim \mathcal{N}(2, 3)$. $\mathbb{P}(-1 \leq X \leq 5)$ vaut :
\begin{answers}
\bad{$\Phi(5) - \Phi(-1)$}
\bad{$\Phi(1) \approx 0{,}841$}
\good{$2\Phi(1) - 1 \approx 0{,}683$}
\bad{$2\Phi\left(\frac13\right) - 1 \approx 0{,}261$}
\end{answers}
\begin{explanations}
On standardise : $Z = \frac{X - 2}{3}$ suit la loi $\mathcal{N}(0,1)$ et
$\mathbb{P}\left(\frac{-1 - 2}{3} \leq Z \leq \frac{5 - 2}{3}\right) = \mathbb{P}(-1 \leq Z \leq 1) = 2\Phi(1) - 1$.

Contre-sens visés : appliquer $\Phi$ aux bornes non standardisées (A) ;
oublier la borne inférieure (B) ; diviser par la variance $9$ au lieu de
l'écart-type $3$, ce qui donne les bornes $\pm\frac13$ (D).
\end{explanations}
\end{question}
}
